% Extended abstract — canonical text from notes/extended_abstract.md (OSF plan draft, WP1)
% Compile with pdflatex (Overleaf-compatible).
\documentclass[11pt,a4paper]{article}

\usepackage[utf8]{inputenc}
\usepackage[T1]{fontenc}
\usepackage[english]{babel}
\usepackage[margin=2.5cm]{geometry}
\usepackage{setspace}
\usepackage[hidelinks]{hyperref}

\title{Digital lifestyles and consumption-based carbon emissions:\\
individual-level evidence from linked bank-transaction, survey, and register data}
\author{David Andersson \and Mathias Lehner% \and G\"oran Finnveden \and Anna Furberg
}
\date{Extended abstract --- \today}

\begin{document}

\maketitle

\onehalfspacing

Whether digitalization lowers or raises the greenhouse gas emissions embedded in
household consumption remains unresolved. Optimists argue that digital time use
substitutes for mobility and material consumption (``dematerialization,
decarbonization, demobilization''); skeptics argue that the time and money freed by
digital substitution rebound into additional consumption, and that digital channels
such as e-commerce actively generate demand, thus spurring more consumption. The
evidence on both sides comes from aggregate panels or self-reported surveys:
existing studies infer travel reduction from time-use diaries, or correlate
broadband penetration with national emissions, but none observes digital behavior
and measured consumption-based emissions in the same individuals. The two
literatures consequently talk past each other, and the net micro-level relationship
--- and, crucially, its composition --- is unknown.

We provide the first individual-level evidence, to our knowledge, linking a
credible measure of digital lifestyle intensity to measured consumption-based
carbon footprints. Approximately 6{,}000 Swedish adults completed surveys in 2024 and
connected their bank accounts to a research app that retrieved transactions back to
2019, categorized them to COICOP, and converted them to CO\textsubscript{2}e. Digital lifestyle
intensity is measured by a pre-registered index of device-use and digital-activity
frequencies from the baseline survey, anchored cardinally by a device-assisted
screen-time report (respondents read average daily screen time from their phone's
own Screen Time / Digital Wellbeing statistics). Register linkage provides income,
education, household composition, and DeSO-level urbanity.

Three results structure the paper. First, the \emph{net gradient}: the association between
digital intensity and total consumption-based CO\textsubscript{2}e, expressed per additional daily
hour of screen time, controlling for sociodemographic factors. Second, the
\emph{decomposition}: whether the gradient is negative in transport (demobilization) and
positive in goods, e-commerce-intensive categories, and digital services (rebound
and direct effects), with no-channel categories (rent, insurance) as placebos ---
it is this pattern, rather than the net sign alone, that adjudicates between the
competing accounts. Third, the \emph{mechanism}: whether lower out-of-home leisure
frequency among the digitally intensive accounts for part of the transport
gradient. Pre-specified heterogeneity by age, gender, and urbanity speaks to where
in the population digitalization is most plausibly emission-reducing.

The results supply the missing empirical foundation for scenario and policy work on
low-carbon digitalization --- within our project, the parameters for prospective
life-cycle scenarios of future digital lifestyles --- and a template for linking
financial-transaction and survey data to study lifestyle transitions at the
individual level.

\end{document}
